\section{Ablation study}\label{sec:ablation}
The ten development queues let us vary the evidence rule while holding the cohort fixed (Table~\ref{tab:ablation}). The permissive rich-claim version connected more fragments, but it produced 2,580 false within-group pairs and 22 mixed groups. Requiring original independent claims eliminated those observed errors in this cohort and reached 158/200 R@20 with the fixed rank bonus. The endpoint fallback then moved R@20 to 169/200 and recoverability from 166/200 to 177/200 on the \emph{same} queues; 14 new unions were genuine under generator truth. It still fell one short of the planned 170/200 gate. None of these counts belongs to the later fresh holdout.
\begin{table}[t]\centering\footnotesize\caption{Matched controlled development ablation (10 queues; 200 top-20 opportunities). All rows use simulator truth; the unsafe variant fails the pair-safety condition.}\label{tab:ablation}
\begin{tabularx}{\textwidth}{@{}Yrrr@{}}\toprule Method & R@20 & Recoverable & False pairs\\\midrule
Rich claims, single lineage & 155/200 & 160/200 & 2,580\\
Independent evidence, rank bonus & 158/200 & 166/200 & 0\\
Independent evidence + endpoint fallback & 169/200 & 177/200 & 0\\\bottomrule
\end{tabularx}\end{table}
Other investigations used separate development cohorts. One hard-negative ranker left groups unchanged but moved matched R@20 from 263/300 to 262/300. Delaying or removing independent evidence also reduced reconstruction. We do not pool those trials with the primary holdout or describe them as holdout ablations. There is still no eligible real-queue comparison of C/P/V against a single LLM, so the contribution of LLM agent count has not been measured.
